\begin{figure}[hbt]
\centering
\caption{Fluxo de organização das tabelas ofensivas e defensivas extraídas do FBref para a construção da base final de atributos dos jogadores.}
\label{fig:data_processing}

\begin{tikzpicture}[
    node distance=1.6cm and 1.8cm,
    box/.style={
        rectangle,
        draw,
        rounded corners,
        align=center,
        minimum width=4.6cm,
        minimum height=1.1cm
    },
    arrow/.style={
        -{Latex[length=3mm]},
        thick
    }
]

% Blocos iniciais
\node (off) [box] {Tabelas de estatísticas\\ofensivas (FBref)};
\node (def) [box, right=of off] {Tabelas de estatísticas\\defensivas (FBref)};

% Bloco central
\node (concat) [box, below=of $(off)!0.5!(def)$] {Concatenação das tabelas};

% Demais etapas
\node (dup) [box, below=of concat] {Remoção de colunas duplicadas};
\node (irr) [box, below=of dup] {Remoção de atributos irrelevantes\\(ex.: Nation)};
\node (final) [box, below=of irr] {Base consolidada por jogador};
\node (csv) [box, below=of final] {Armazenamento em arquivos CSV por time e temporada\\(ex.: \textit{data\_napoli\_2021-2022.csv})};

% Setas
\draw [arrow] (off.south) -- (concat.north west);
\draw [arrow] (def.south) -- (concat.north east);
\draw [arrow] (concat) -- (dup);
\draw [arrow] (dup) -- (irr);
\draw [arrow] (irr) -- (final);
\draw [arrow] (final) -- (csv);

\end{tikzpicture}

\fonte{Elaborado pelo autor.}
\end{figure}